% Capitulo: Com que frequência os dados precisam atualizar?
\chapter{Com que frequência os dados precisam atualizar?}
\label{cap:frequencia}

\rascunho{
  \textbf{Ideia do capítulo.} Uma das decisões mais importantes de um produto de dados,
  e uma das mais tomadas no automático: com que frequência ele atualiza. A regra que
  quero defender é atualizar o mais rápido que as fontes permitem, mas só o quanto a
  decisão de negócio precisa. Um relatório que atualiza de hora em hora e é aberto uma
  vez por semana gasta recurso sem gerar valor.
  \begin{itemize}
    \item Exemplo de abertura, da minha experiência como gerente: um pedido de
      \enquote{tempo real} que, olhando o uso, era uma reunião semanal. \pendente{escolher o caso}
    \item A frequência é uma definição de negócio antes de ser técnica, e pode estar
      errada. Por isso os produtos entregues precisam de métricas sobre eles mesmos
      (uso, atualizações, custo), e a definição precisa ser revisada.
    \item Na era da IA, em que quase tudo é cobrado em créditos e tokens, errar a
      frequência custa dinheiro de forma direta (ver o capítulo~\ref{cap:custo-consumo}).
    \item Conceito emprestado de SRE: objetivo de nível de serviço (SLO), aqui aplicado
      à atualização do dado \parencite{beyer2016}. Atualização (\emph{freshness}) como
      uma das dimensões da observabilidade de dados \parencite{moses2022}.
  \end{itemize}
}

% ---------------------------------------------------------------------------
\section{O ritmo da decisão}
\label{sec:ritmo-decisao}

\rascunho{
  \begin{itemize}
    \item A pergunta certa para o stakeholder não é \enquote{com que frequência você
      quer que atualize?} (a resposta é sempre \enquote{o mais rápido possível}), mas
      \enquote{quando e como você toma a decisão com esse número?}.
    \item Decisões operacionais (despacho, estoque do dia, fraude) pedem atualização
      frequente; decisões táticas e estratégicas (metas, orçamento, planejamento de
      safra) convivem bem com atualização diária ou semanal.
    \item A Figura~\ref{fig:frequencia-consumo} cruza a frequência de atualização com a
      frequência de consumo: os dois quadrantes desalinhados são desperdício e risco.
    \item Mostrar no próprio dashboard a data e a hora da última atualização: o usuário
      precisa saber a idade do número que está vendo.
  \end{itemize}
}

\pendente{escrever}

\begin{figure}[htbp]
  \centering
  % Frequencia de atualizacao contra frequencia de consumo (cap. 16).
\begin{tikzpicture}[diagrama, x=1mm, y=1mm]
  % quadrantes (os da diagonal, alinhados, em branco; os desalinhados, em cinza)
  \node[caixa, fill=fundo, minimum width=46mm, minimum height=26mm, text width=40mm] at (25,43)
    {\textbf{Desperdício}\\[1mm]atualiza de hora em hora, alguém olha uma vez por semana};
  \node[caixa, minimum width=46mm, minimum height=26mm, text width=40mm] at (75,43)
    {\textbf{Operacional}\\[1mm]a decisão é de minutos ou horas; atualizar rápido vale o custo};
  \node[caixa, minimum width=46mm, minimum height=26mm, text width=40mm] at (25,15)
    {\textbf{Estratégico}\\[1mm]decisão mensal ou semanal; atualização periódica basta};
  \node[caixa, fill=fundo, minimum width=46mm, minimum height=26mm, text width=40mm] at (75,15)
    {\textbf{Risco}\\[1mm]muita gente decide todo dia com um número velho};

  % eixos
  \draw[seta] (-3,-1) -- (104,-1);
  \draw[seta] (-1,-3) -- (-1,61);
  \node[rotulo, anchor=north] at (50,-2) {com que frequência alguém olha os números};
  \node[rotulo, anchor=south, rotate=90] at (-2,29) {com que frequência os dados atualizam};
  \node[rotulo, anchor=north west] at (0,-2) {menos};
  \node[rotulo, anchor=north east] at (102,-2) {mais};
\end{tikzpicture}

  \caption{Frequência de atualização contra frequência de consumo. Nos quadrantes em
    branco as duas estão alinhadas; nos quadrantes em cinza, não.}
  \label{fig:frequencia-consumo}
  \fonte{elaborado pelo autor.}
\end{figure}

% ---------------------------------------------------------------------------
\section{O que as fontes permitem}
\label{sec:limites-fontes}

\rascunho{
  \begin{itemize}
    \item O teto da frequência é dado pelas fontes: janelas de extração, limites de API,
      sistemas que só fecham o dado no fim do dia, captura de mudanças (CDC) disponível
      ou não.
    \item Lote ou streaming: o custo e a complexidade crescem muito quando se passa de
      diário para horário e de horário para contínuo \parencite{reis2022}.
    \item Carga incremental contra carga completa: a mesma frequência pode custar dez
      vezes menos com um modelo incremental. \pendente{conferir se uso essa ordem de grandeza ou um exemplo próprio}
    \item Atualizar mais rápido que a fonte muda não traz dado novo, só custo.
  \end{itemize}
}

\pendente{escrever}

% ---------------------------------------------------------------------------
\section{O que cada profissional precisa olhar}
\label{sec:papeis-frequencia}

\rascunho{
  A Tabela~\ref{tab:papeis-frequencia} é um rascunho para revisar com a minha
  experiência. A ideia é que a frequência não é responsabilidade de um papel só: o
  gerente decide, mas cada papel enxerga uma parte da conta.
}

\pendente{escrever}

\begin{table}[htbp]
  \centering
  \caption{A frequência de atualização vista por cada papel.}
  \label{tab:papeis-frequencia}
  \small
  \begin{tabularx}{\textwidth}{@{}>{\raggedright\arraybackslash}p{0.27\textwidth}
      >{\raggedright\arraybackslash}X@{}}
    \toprule
    Papel & Com o que precisa se preocupar \\
    \midrule
    Gerente ou líder de dados & Entender a necessidade real do negócio, negociar a
      frequência com o stakeholder, equilibrar custo e valor, acompanhar as métricas de
      uso e revisar a decisão de tempos em tempos. \\
    \addlinespace
    Stakeholder (dono do negócio) & Dizer quando e como decide com o número, e aceitar
      que atualizar mais rápido tem custo. \\
    \addlinespace
    Analista ou desenvolvedor de BI & Desenhar o produto para o ritmo de quem consome,
      mostrar a data da última atualização e não prometer uma frequência que a fonte não
      sustenta. \\
    \addlinespace
    Analytics engineer & Escolher a materialização certa (visão, tabela, incremental)
      para a frequência combinada e testar a atualização do dado. \\
    \addlinespace
    Engenheiro de dados & Conhecer os limites das fontes, o custo de cada execução e
      monitorar atrasos e falhas de carga. \\
    \bottomrule
  \end{tabularx}
  \fonte{elaborado pelo autor.}
\end{table}

% ---------------------------------------------------------------------------
\section{Métricas sobre os próprios produtos de dados}
\label{sec:metricas-uso}

\rascunho{
  \begin{itemize}
    \item Tratar dashboards, tabelas e modelos como produtos, com métricas próprias:
      acessos, usuários distintos, data do último acesso, número de atualizações, custo
      por execução, custo por acesso.
    \item A métrica que mais revela desperdício: atualizações por acesso, ou créditos
      por acesso. A consulta abaixo mostra a ideia, sobre tabelas de metadados que quase
      toda ferramenta de BI e armazém de dados oferece de alguma forma.
      \pendente{trocar os nomes pelos metadados da ferramenta usada como exemplo}
    \item Produtos sem acesso há semanas são candidatos a desligar, ou a baixar a
      frequência.
  \end{itemize}
}

\pendente{escrever}

\begin{sql}
-- Atualizações contra acessos de cada produto de dados nos últimos 30 dias
WITH atualizacoes AS (
    SELECT produto,
           COUNT(*)            AS qtd_atualizacoes,
           SUM(custo_creditos) AS creditos
    FROM meta_atualizacoes
    WHERE executado_em >= CURRENT_DATE - INTERVAL '30' DAY
    GROUP BY produto
),
acessos AS (
    SELECT produto,
           COUNT(*)                AS qtd_acessos,
           COUNT(DISTINCT usuario) AS usuarios
    FROM meta_acessos
    WHERE acessado_em >= CURRENT_DATE - INTERVAL '30' DAY
    GROUP BY produto
)
SELECT a.produto,
       a.qtd_atualizacoes,
       COALESCE(v.qtd_acessos, 0)            AS qtd_acessos,
       COALESCE(v.usuarios, 0)               AS usuarios,
       a.creditos,
       a.creditos / NULLIF(v.qtd_acessos, 0) AS creditos_por_acesso
FROM atualizacoes AS a
LEFT JOIN acessos AS v ON v.produto = a.produto
ORDER BY creditos_por_acesso DESC NULLS FIRST;  -- sem acesso nenhum aparece primeiro
\end{sql}

% ---------------------------------------------------------------------------
\section{Revisar a definição}
\label{sec:revisar-frequencia}

\rascunho{
  \begin{itemize}
    \item A frequência combinada no começo do projeto é uma hipótese. As métricas de uso
      dizem se ela estava certa.
    \item Ritual simples: revisar mensal ou trimestralmente os produtos mais caros e os
      menos acessados. Encaixa na revisão do Scrum com Kanban (ver o
      capítulo~\ref{cap:scrumban}).
    \item Baixar a frequência ou desligar um produto também é entrega, e precisa ser
      comunicado ao stakeholder com os números de uso na mão.
    \item Relato: uma revisão dessas e o que ela economizou. \pendente{escolher o caso; números só em ordem de grandeza}
  \end{itemize}
}

\pendente{escrever}

% ---------------------------------------------------------------------------
\begin{leituras}
  \leitura{beyer2016}{De onde vem a ideia de objetivo de nível de serviço (SLO). Os
    capítulos sobre SLOs se aplicam quase sem adaptação a prazos de atualização de dados.}
  \leitura{moses2022}{Trata a atualização do dado como uma das dimensões da
    observabilidade e mostra como monitorar atrasos antes do usuário reclamar.}
  \leitura{reis2022}{Para a discussão de lote contra streaming e do custo de cada escolha
    no ciclo de vida da engenharia de dados.}
\end{leituras}
